\section{Problem Summary}
Determine how many zeros the decimal representation of \(n!\) ends with, without ever constructing the factorial explicitly.

\section{Editorial}
A trailing zero comes from a factor of \(10\), which is \(2 \times 5\).

In \(n!\), factors of \(2\) are much more common than factors of \(5\), so the answer is exactly the number of times \(5\) appears in the prime factorization of \(n!\).

Count those \(5\)-factors in layers:
\begin{itemize}
\item every multiple of \(5\) contributes at least one factor of \(5\)
\item every multiple of \(25\) contributes one extra factor
\item every multiple of \(125\) contributes yet another extra factor
\item and so on
\end{itemize}

So the answer is
\[
\left\lfloor \frac{n}{5} \right\rfloor +
\left\lfloor \frac{n}{25} \right\rfloor +
\left\lfloor \frac{n}{125} \right\rfloor + \cdots
\]
until the divisor exceeds \(n\).

\section{Complexity Analysis}
\begin{itemize}
\item Time: \(O(\log_5 n)\)
\item Memory: \(O(1)\)
\end{itemize}

\section{Notes}
This is a counting problem, not a big-integer problem. Never try to build \(n!\) directly.
